\section{Задача~5c}

При каких $m$ и $n$ поле $\FFpm$ содержит $\FFpn$?

\begin{solution}
	Вспомним, что такое $\FFpm$. При $m = 1$
	\[
		\FFp := \ZZ/(p),
	\]
	а при $m > 1$ можно положить по определению
	\[
		\FFpm := \frac{\FFp{}[t]}{(f_m)} = \FFp{}[\theta]
		,
	\]
	где $f_m \in \FFp{}[t]$~"--- неприводимый многочлен степени~$m$ над полем~$\FFp$, а $f(\theta) = 0$.

	% Пусть, к примеру,
	% \[
	% 	f_m(t) =
	% 	b_0 + b_1 t + \dots + b_{m-1}t^{m-1} + b_m t^m
	% 	.
	% \]
	% Тогда если~$f_m(\theta) = 0$, то
	% \[
	% 	b_0 + b_1 \theta + \do\thetas + b_{m-1}\theta^{m-1} + b_m \theta^m = 0,
	% \]
	% откуда
	% \[
	% 	\theta^m =
	% 	-
	% 	\frac{1}{b_m}
	% 	\,
	% 	\bigl(
	% 		b_0 + b_1 \theta + \dots + b_{m-1}\theta^{m-1}
	% 	\bigr)
	% 	.
	% \]
	% Степень~$\theta^m$ выражается через предыдущие.

	% Поэтому произвольный элемент~$a \in \FFpm$ имеет вид
	% \[
	% 	a = a_0 + a_1 \theta + \dots + a_{m-1}\theta^{m-1},
	% \]
	% причём~$\theta^m$ выражается через степени $\theta^\mu$ при $\mu < m$, как указано в выражении~$f_m(\theta) = 0$.

	Докажем, что
	\[
		(\FFpn \subset \FFpm)
		\iff
		(n \mid m)
	\]
	(здесь <<$\mid$>> --- знак, означающий <<делит>>)%
	\footnote{%
		Говорят, что $a \in K$ делит~$b \in K$, если $b = ka$ для некоторого~$k \in K$.
	}.

	Пусть $\FFpn \subset \FFpm$. Тогда
	\[
		[\FFpm : \FFp] =
		[\FFpm : \FFpn] \cdot
		[\FFpn : \FFp]
		.
	\]
	Так как $[\FFpm : \FFp] = m$ и $[\FFpn : \FFp] = n$, то
	\[
		m = [\FFpm : \FFpn] \cdot n
		,
	\]
	то есть $n \mid m$.

	Обратно, пусть $n \mid m$, то есть $m = kn$. Тогда для любого элемента~$a \in \FFpn$
	\[
		a^{p^m} =
		a^{p^{kn}} =
		a^{(p^n)^k} =
		a^{p^n \cdot (p^n)^{(k-1)}} =
		a^{(p^n)^{(k-1)}} =
		\dots
		a^{p^n} =
		a,
	\]
	так что~$a \in \FFpm$. Поэтому каждый корень~$f_n(x) := (x^{p^n} - x)$ является корнем~$f_m(x) := (x^{p^m} - x)$, а значит, в разложении на линейные множители каждый множитель, входящий в состав~$f_n$, входит и в состав~$f_m$, так что $f_n \mid f_m$.
	% \[
	% 	\bigl(
	% 		x^{p^n} - x
	% 	\bigr)
	% 	\mid
	% 	\bigl(
	% 		x^{p^m} - x
	% 	\bigr)
	% 	,
	% \]
	% так как 
\end{solution}
